\section{Passage to arbitrary von Neumann algebras}
\label{gen:section}

We first remove the separability assumption from the state-preserving
result. We then treat the part of an inclusion on which no normal
positive functional centralizes the MASA. The two arguments will be
combined using corners by projections in the MASA; these projections
need not be central in the ambient algebra. The support $e$ of the
centralizing functionals will carry an expected MASA, while its complement
$h$ has none of those functionals. Only after choosing the finite vector
tests will we split $e=f+g$, with a faithful state on the $f$-corner
and a residual projection $g$ that annihilates the test vectors.

\subsection{Separable reduction}

\begin{lemma}\label{lem:separable-reduction}
The conclusion of Theorem~\ref{thm:state-paving} holds without the
assumption of separable predual, with the same partition-size bound.
\end{lemma}

\begin{proof}
Let $\varphi$ be a faithful normal state on $M$, let $A\subset M$ be a
MASA in its centralizer, and fix a self-adjoint contraction $x\in M$.
Construct countable sets of generators as follows. Start with $\one,x$.
At each stage form the unital $\mathbb Q(i)$-algebra generated by the
current set and its adjoints. For every element $t$ of this countable
algebra, add $E_A(t)$ and the projections of finite $A$-partitions
$\mathcal Q_{t,k}$ satisfying
\[
 \norm{C_{\mathcal Q_{t,k}}(t)-E_A(t)}_\varphi^\#<1/k,
 \qquad k\geq1.
\]
These partitions exist by Lemma~\ref{lem:compression-net}. Repeat this
procedure countably many times, let $B$ be the union of the resulting
countable algebras, and set
\[
 M_1=W^*(B),\qquad A_1=A\cap M_1.
\]

Every element of $M_1$ is the limit of an operator-norm bounded sequence
from $B$ in $\norm{\cdot}_\varphi^\#$. Indeed, Kaplansky density gives
bounded strong-star approximants from the norm closure of $B$, and one
can then approximate in norm by elements of $B$. On bounded sets the
state sharp seminorm determines the strong-star topology, so sequences
suffice here. Taking real and imaginary parts separately gives this
assertion for arbitrary elements. Since $E_A$ is contractive for the
state sharp seminorm and in operator norm, it follows that
$E_A(M_1)\subset A_1$.

Suppose $z\in M_1$ commutes with $A_1$. Choose bounded $t_j\in B$
converging to $z$ in the state sharp seminorm, and choose one of the
stored partitions $\mathcal Q_j\subset A_1$ for each $t_j$ so that its
compression error is less than $1/j$. Since $C_{\mathcal Q_j}(z)=z$,
contractivity gives
\[
 \norm{z-E_A(z)}_\varphi^\#
 \leq 2\norm{z-t_j}_\varphi^\#
       +\norm{C_{\mathcal Q_j}(t_j)-E_A(t_j)}_\varphi^\#
 \longrightarrow0.
\]
Thus $z=E_A(z)\in A_1$, proving that $A_1$ is a MASA in $M_1$.
It lies in the centralizer of $\varphi|_{M_1}$. The GNS representation
of this restricted faithful normal state is faithful and normal, and
its Hilbert space is separable because the countable set $B$ applied
to the cyclic vector is dense. Hence $M_1$ has separable predual.
The restriction of $E_A$ is its state-preserving expectation onto
$A_1$. Applying Theorem~\ref{thm:state-paving} in $M_1$ proves the
assertion, including the unchanged bound on the partition size.
\end{proof}

\subsection{The corner without centralizing normal functionals}

We use standard form in the next argument. Write $a\xi b$ for the
commuting left and right actions, with the right action of $b$ given
by $Jb^*J$. Every normal positive functional $\theta$ has a unique
representative $\xi_\theta$ in the natural positive cone. This cone
is selfdual and is preserved by $\xi\mapsto u\xi u^*$ for a unitary
$u$. Moreover,
\begin{equation}\label{gen:powers-stormer}
 \norm{\xi_\theta-\xi_\psi}^{2}
 \leq\norm{\theta-\psi}_{M_*}.
\end{equation}
These standard-form facts hold for arbitrary von Neumann algebras;
see \cite[Theorem~3.12]{Hiai20}. The inequality also appears in
\cite[Theorem~4(8), Equation~(5.10)]{Araki74}. For two functionals in
a non-$\sigma$-finite algebra, it can also be applied in the corner
supported by $\supp(\theta+\psi)$, where their sum is faithful.

\begin{lemma}\label{lem:no-central-functional}
Let $D\subset P$ be a MASA such that no nonzero normal positive
functional on $P$ centralizes $D$. For every normal state $\theta$ on
$P$, every integer $m\geq2$, and every $\alpha>0$, there are a positive
contraction $s\in P$ and an $m$-partition $(p_i)_{i=1}^m$ in $D$ such
that
\begin{equation}\label{gen:small-diagonal}
 \theta(s)>1-\alpha,\qquad p_i s p_i\leq m^{-1}p_i
 \quad(1\leq i\leq m).
\end{equation}
\end{lemma}

\begin{proof}
In standard form, each finite partition $\mathcal Q$ in $D$ defines
an orthogonal projection
\[
 P_{\mathcal Q}\xi=\sum_{q\in\mathcal Q}q\xi q.
\]
These projections decrease under refinement to the projection onto
the $D$-central vectors. A nonzero central vector $\zeta$ would give
a nonzero $D$-central normal positive functional
$a\mapsto\ip{a\zeta}{\zeta}$: conjugation by a unitary of $D$ can be
transferred to the commuting right action. The hypothesis therefore
implies $P_{\mathcal Q}\to0$ strongly along the net of finite
partitions.

Choose independent random labels $Z_q$, uniform on the $m$th roots
of unity, and put $u=\sum_{q\in\mathcal Q}Z_q q$. For $1\leq k<m$,
independence and $\E Z_q^k=0$ give
\[
 \E\ip{u^k\xi_\theta u^{-k}}{\xi_\theta}
   =\ip{P_{\mathcal Q}\xi_\theta}{\xi_\theta}
   \longrightarrow0.
\]
Each overlap is real and nonnegative, by the natural-cone properties.
Consequently, for a sufficiently fine partition there is one choice
of labels making all these overlaps arbitrarily small. The states
\[
 \theta_j=\theta\circ\Ad(u^{-j}),\qquad 0\leq j<m,
\]
have natural-cone vectors $\xi_j=u^j\xi_\theta u^{-j}$. Their pairwise
overlaps are among those just considered. Equation~\eqref{gen:powers-stormer}
shows that $\norm{\theta_j-\theta_k}_{P_*}$ is arbitrarily close to
$2$ for every $j\ne k$.

Fix a small $\eta>0$. The positive support projection of
$\theta_j-\theta_k$ attains half the norm of that Hermitian
functional, since its value at $\one$ is zero. Thus the labels can
be chosen so that there are projections $q_{jk}\in P$ with
\begin{equation}\label{gen:pair-separation}
 \theta_j(q_{jk})\geq1-\eta,\qquad
 \theta_k(q_{jk})\leq\eta\qquad(j\ne k).
\end{equation}
For $0<t<1$ put
\[
 a_j=\frac1{m-1}\sum_{k\ne j}q_{jk},\qquad
 b_j=1_{[1-t,1]}(a_j).
\]
Spectral calculus and Equation~\eqref{gen:pair-separation} yield
\begin{equation}\label{gen:spectral-separation}
 \theta_j(b_j)\geq1-\frac\eta t,\qquad
 \theta_k(b_j)^{1/2}\leq\sqrt\eta+\sqrt{(m-1)t}\quad(k\ne j).
\end{equation}
For the first estimate, use $\theta_j(1-a_j)\leq\eta$ and
$1-b_j\leq t^{-1}(1-a_j)$. For the second, the positive-summand
inequality $1-q_{jk}\leq(m-1)(1-a_j)$ gives
\[
 \norm{(1-q_{jk})b_j}^2
   =\norm{b_j(1-q_{jk})b_j}\leq(m-1)t.
\]
Splitting $b_j\xi_k=b_jq_{jk}\xi_k+b_j(1-q_{jk})\xi_k$ proves
the claim, using the adjoint of this operator-norm estimate.

The pairwise separating projections need not be mutually orthogonal,
so they cannot simply be added to form a partition. Sequential products
give positive effects with sum at most the identity while retaining
almost full mass for their respective states. Define
\[
 V_0=\one,\qquad V_j=(1-b_{j-1})\cdots(1-b_0),\qquad
 s_j=V_j^*b_jV_j\quad(0\leq j<m).
\]
These positive elements need not be projections. They satisfy
\begin{equation}\label{gen:effects}
 s_j=V_j^*V_j-V_{j+1}^*V_{j+1},\qquad
 \sum_{j=0}^{m-1}s_j=\one-V_m^*V_m\leq\one.
\end{equation}
Telescoping the products of contractions also gives
\[
 \norm{V_j\xi_j-\xi_j}
 \leq\sum_{k<j}\norm{b_k\xi_j}.
\]
By Equation~\eqref{gen:spectral-separation}, whenever $\eta<t$,
\[
 \norm{b_jV_j\xi_j}
 \geq\sqrt{1-\eta/t}
       -j\bigl(\sqrt\eta+\sqrt{(m-1)t}\bigr).
\]
For fixed $m$, choosing $t$ small and then $\eta$ sufficiently small
makes every $\theta_j(s_j)=\norm{b_jV_j\xi_j}^2$ arbitrarily close
to $1$.

Set
\[
 s=\frac1m\sum_{j=0}^{m-1}u^{-j}s_ju^j.
\]
Then $0\leq s\leq\one$ and $\theta(s)$ is the average of the
$\theta_j(s_j)$, so it exceeds $1-\alpha$ with the choices above.
Since $u^m=\one$, Equation~\eqref{gen:effects} gives
\begin{equation}\label{gen:cyclic-average}
 \sum_{\ell=0}^{m-1}u^\ell s u^{-\ell}
 =\frac1m\sum_{k=0}^{m-1}u^k
       \left(\sum_{j=0}^{m-1}s_j\right)u^{-k}
 \leq\one.
\end{equation}
Let $p_1,\ldots,p_m$ be the spectral projections of $u$ for the
$m$th roots, allowing zero projections. Compressing
Equation~\eqref{gen:cyclic-average} by $p_i$ gives
$mp_i s p_i\leq p_i$, as required.
\end{proof}

\begin{corollary}\label{gen:no-central-approximation}
Under the hypotheses of Lemma~\ref{lem:no-central-functional}, let
$P$ act normally on a Hilbert space and let $x=x^*\in P$ satisfy
$\norm{x}\leq1$. For each $m\geq2$, the operators $sxs$ supplied
by that lemma approximate $x$ strongly, with associated
$m$-partitions $\mathcal P$ satisfying
\[
 \norm{sxs}\leq1,\qquad
 \norm{C_{\mathcal P}(sxs)}\leq1/m.
\]
Here the approximation is on any prescribed finite set of vectors,
to any positive accuracy, while $m$ stays fixed.
\end{corollary}

\begin{proof}
For test vectors $\xi_1,\ldots,\xi_d$, normalize the sum of their
vector functionals and those of $x\xi_1,\ldots,x\xi_d$ to obtain
a normal state $\theta$. If this sum vanishes, all tests are zero
and one may take $s=0$. Otherwise Lemma~\ref{lem:no-central-functional}
makes
\[
 \theta((1-s)^2)\leq1-\theta(s)
\]
arbitrarily small. Thus $1-s$ is arbitrarily small on every test
vector and its image under $x$. The identity
\[
 sxs-x=sx(s-1)+(s-1)x
\]
proves the strong approximation. Finally,
$-s^2\leq sxs\leq s^2$ and $s^2\leq s$ imply
$-p_i/m\leq p_i sxs p_i\leq p_i/m$. The norm of the block diagonal
sum is the maximum of its block norms.
\end{proof}

\subsection{Assembly of the three corners}

\begin{lemma}\label{gen:support-decomposition}
For any MASA $A\subset M$, set
\[
 e=\bigvee\{\supp(\lambda):\lambda\in M_*^+,
                  \ \lambda\text{ centralizes }A\},\qquad h=\one-e.
\]
Then $e,h\in A$. If $e\ne0$, there is a faithful normal conditional
expectation $E:eMe\to Ae$. If $h\ne0$, the inclusion
$Ah\subset hMh$ has no nonzero centralizing normal positive functional.
\end{lemma}

\begin{proof}
The support of an $A$-central normal positive functional is invariant
under every unitary in $A$, and hence belongs to $A$ by maximal
abelianness. Such a functional is faithful on its support corner,
and remains faithful and centralizing upon restriction to any nonzero
$A$-subprojection of its support.

Choose a maximal orthogonal family $(q_\beta)$ of nonzero projections
in $A$, each subordinate to the support of one such functional.
Their sum is $e$: a nonzero remainder under $e$ would have nonzero
intersection with one of the defining supports, contradicting
maximality. On each $q_\beta Mq_\beta$, normalize the restricted
functional to a faithful normal state. The corner $Aq_\beta$ is a
MASA in its centralizer. Takesaki's conditional expectation theorem
\cite{Takesaki72}, in the state formulation
\cite[Theorem~6.6]{Hiai20}, gives a faithful normal expectation
$E_\beta:q_\beta Mq_\beta\to Aq_\beta$. Its hypotheses hold because
the centralizer is fixed by the modular group and a state is finite
on the subalgebra.

Define
\begin{equation}\label{gen:glued-expectation}
 E(z)=\sum_\beta E_\beta(q_\beta zq_\beta),\qquad z\in eMe.
\end{equation}
The orthogonal block sum is well defined, bounded in norm by
$\norm{z}$, and belongs to $Ae$. It is positive, unital, $Ae$-bimodular,
and fixes $Ae$. For a bounded increasing net of positive elements,
normality of every $E_\beta$ identifies the supremum of every
$q_\beta$-block; this proves normality of $E$, even for an uncountable
family. If $z\geq0$ and $E(z)=0$, faithfulness on each corner gives
$q_\beta zq_\beta=0$, hence $z^{1/2}q_\beta=0$ for all $\beta$.
Since $\sum_\beta q_\beta=e$, this forces $z=0$. Thus $E$ is faithful.

An $Ah$-central normal positive functional on $hMh$ would extend by
$z\mapsto hzh$ to an $A$-central one on $M$ supported under $h$.
The definition of $e$ excludes a nonzero such functional. All corner
MASAs used here remain maximal abelian: an element in a corner
commuting with the corner of $A$, extended by zero, commutes with $A$.
\end{proof}

\begin{proof}[Proof of Theorem~\ref{thm:main}]
If $x=0$, take $y=a=0$ and the one-part partition, padded with zero
projections to the size chosen below. Otherwise, by
scaling it suffices to prove the assertion for $\norm{x}=1$, with an
arbitrary strong-approximation tolerance $\delta>0$. Fix
$0<\varepsilon<1$ and put $\gamma=\varepsilon/4$. Let $r_f$ be the
partition-size bound $R_{\mathrm{st}}(\gamma)$ in
Theorem~\ref{thm:state-paving}, and fix $m\geq2$ with
$1/m\leq\gamma$. We shall use
\begin{equation}\label{gen:final-size}
 R(\varepsilon)=r_f+m+1.
\end{equation}
More explicitly, the choices in the state-paving proof may be taken as
\[
 b=\sigma=\frac{\varepsilon}{96(K+1)},\qquad
 r=\left\lceil\frac{32}{b^2}\right\rceil,\qquad
 n=\left\lceil\frac{1024}{\varepsilon^2}\right\rceil,\qquad
 m=\left\lceil\frac4\varepsilon\right\rceil.
\]
Here $2\sqrt{2/r}+2/r\leq b/2+b^2/16<b$,
$4/\sqrt n\leq\gamma/2$, and $1/m\leq\gamma$.
Consequently $r_f=(r+1)^2n$ and
\begin{equation}\label{gen:sixth-power}
 R(\varepsilon)=(r+1)^2n+m+1
       \leq C_{\mathrm{ap}}\varepsilon^{-6}
       \qquad(0<\varepsilon<1)
\end{equation}
for a universal constant $C_{\mathrm{ap}}$. Only the universal buffer
$K$ enters this constant; graph, slab, and approximation parameters
do not enter the color count.
These choices precede the finite test set and its tolerance.

Let $e,h$ and $E$ be as in Lemma~\ref{gen:support-decomposition}.
Given a finite set $F$ of test vectors, choose a unit vector $\xi_0$
such that $\norm{x\xi_0}>3/4$, and put $F'=F\cup\{\xi_0\}$.
On $Ae$ consider the normal positive functional
\[
 \psi(a)=\sum_{\xi\in F'}\ip{a e\xi}{e\xi},\qquad a\in Ae,
\]
and let $f=\supp(\psi)\in Ae$ and $g=e-f$. The equality $\psi(g)=0$
implies $g\xi=0$ for every $\xi\in F'$. Empty corners below are
interpreted as zero.

If $f\ne0$, then $\psi|_{Af}$ is faithful. After normalizing it,
compose it with the restriction $E:fMf\to Af$ to obtain a faithful
normal state $\varphi_f$ on $fMf$. Bimodularity of $E$ and
commutativity of $Af$ show that $Af$ is in the
$\varphi_f$-centralizer. Apply Theorem~\ref{thm:state-paving} and
Lemma~\ref{lem:separable-reduction} to the self-adjoint contraction
$x_f=fxf$. They give
\[
 y_f=x_f-w,\qquad w=w^*,\quad\norm{w}\leq2,
 \qquad
 \norm{C_{\mathcal P_f}(y_f)-d_f}\leq\gamma,
 \quad d_f=E(x_f),
\]
where $\mathcal P_f$ is an $r_f$-partition of $f$, and $w$ can be
chosen arbitrarily small on the vectors $f\xi$ for $\xi\in F'$.
Indeed, the bounded self-adjoint corrections in that theorem have
state norm tending to zero, hence converge strongly in the given
normal representation.

On $g$ put $y_g=0$. On $h$ put $x_h=hxh$ and use
Corollary~\ref{gen:no-central-approximation} to choose
$y_h=sx_hs$ arbitrarily close to $x_h$ on the vectors $h\xi$,
$\xi\in F'$, together with an $m$-partition $\mathcal P_h$ of $h$
such that
\[
 \norm{y_h}\leq1,\qquad\norm{C_{\mathcal P_h}(y_h)}\leq\gamma.
\]

Keep all off-diagonal blocks of $x$ between $f,g,h$, and replace its
diagonal blocks by the operators just constructed. Explicitly, set
\[
 y=y_f+y_g+y_h+
       \sum_{\substack{b,c\in\{f,g,h\}\\b\ne c}}bxc.
\]
This operator is self-adjoint. The difference $y-x$ is block diagonal:
its $f$-block has norm at most $2$, its $g$-block has norm at most $1$,
and its $h$-block has norm at most $2$. Hence
\begin{equation}\label{gen:norm-control}
 \norm{y-x}\leq2,\qquad\norm{y}\leq3.
\end{equation}
The $g$-block of the difference annihilates $F'$. Choosing the other
two approximation errors sufficiently small therefore ensures
\begin{equation}\label{gen:test-control}
 \norm{(y-x)\xi}<\min\{\delta,1/4\}\qquad(\xi\in F').
\end{equation}
In particular $\norm{y}>1/2$, by applying this inequality to
$\xi_0$.

Use the projections of $\mathcal P_f$, the projections of
$\mathcal P_h$, and the single projection $g$ as distinct members
of one $A$-partition $\mathcal P$, padding with zeros in absent
corners. It has the size in Equation~\eqref{gen:final-size} and kills
every off-diagonal block between $f,g,h$. With $d=d_f$ on $f$ and
zero on $g+h$, we obtain
\[
 d=d^*\in A,\qquad\norm{C_{\mathcal P}(y)-d}\leq\gamma.
\]
Since compression is norm-contractive, $\norm{d}\leq\norm{y}+\gamma$.
Clip $d$ by functional calculus to the interval
$[-\norm{y},\norm{y}]$, obtaining $a\in A$ with
$\norm{a}\leq\norm{y}$ and $\norm{a-d}\leq\gamma$. It follows that
\[
 \norm{C_{\mathcal P}(y)-a}
 \leq2\gamma=\varepsilon/2\leq\varepsilon\norm{y}.
\]
Equations~\eqref{gen:norm-control} and~\eqref{gen:test-control}
give the remaining requirements. The partition size depends only on
$\varepsilon$, since the bound in Theorem~\ref{thm:state-paving} is
uniform over the state-preserving inclusions used here. Finally,
scaling $y$ and $a$ by the original $\norm{x}$, and first dividing
the requested strong tolerance by that norm, proves the theorem with
the universal constant $\kappa=3$.
\end{proof}

\subsection{Consequences for related paving formulations}
\label{gen:consequences}

The main theorem concerns norm paving of bounded approximants. The two
introductory consequences use the normal-expectation hypothesis to pass
between that formulation and strong-operator paving of the original
operator. The quadratic consequence additionally uses the separate
companion theorem.

\begin{proof}[Proof of Corollary~\ref{cor:so-paving}]
By Theorem~\ref{thm:main}, every such \(x\) is
\((\varepsilon/12,R(\varepsilon/12);3)\)-approximately pavable
in the terminology of \cite[Definition~2.3]{PV15}.
Apply \cite[Proposition~2.4(1)]{PV15} with
\(\varepsilon'=\varepsilon/6\).
Its error is \(2\cdot3\cdot(\varepsilon/6)=\varepsilon\), so
\(R_s(\varepsilon)=R(\varepsilon/12)\) suffices.
\end{proof}

\begin{proof}[Proof of Corollary~\ref{cor:quadratic-expected-approximation}]
By \cite[Theorem~1.1]{OAQuadratic26}, every self-adjoint \(x\in M\)
is \((\varepsilon/2,r_{\varepsilon/2})\) so-pavable, where
\(r_{\varepsilon/2}\leq C\varepsilon^{-2}\) for a universal \(C\).
Since \(\varepsilon/2<\varepsilon\),
\cite[Proposition~2.4(2)]{PV15} makes \(x\)
\((\varepsilon,r_{\varepsilon/2};3)\)-approximately pavable.
Its proof supplies self-adjoint approximants and diagonal witnesses,
with paving error relative to the approximant's norm.
The strong-limit net in \cite[Definition~2.3]{PV15} supplies the
finite-vector conclusion with the fixed count
\(N(\varepsilon)=r_{\varepsilon/2}\).
\end{proof}
